\chapter{Типизированное вложение спектральной самоэнергии в 43-мерную ячейку}
\label{ch:v10-inflow-spectral-self-energy-k43-typed-embedding}

\section{Проверяемая гипотеза}

Предыдущий гейт построил абстрактный двухмодовый резервуар с ненулевой
интенсивной самоэнергией, но не связал его с носителем программы. Настоящий
гейт проверяет, содержится ли такая пара уже в столкновительной ячейке
\begin{equation}
 \mathcal K_{43}=\mathbb C|0\rangle\oplus
 \operatorname{span}\{|a\rangle:a=1,\ldots,42\},
 \label{eq:v10-k43-cell}
\end{equation}
построенной в Томе VIII, и совместима ли она с калибровочной типизацией,
звёздным взаимодействием и шестимерным KMS-множителем.

\section{Каноническая двумерная подячейка}

Полный вещественный шумовой репер состоит из тридцати направлений переноса и
двенадцати калибровочных направлений. Последнее направление есть уже
унаследованный генератор гиперзаряда $Y$. Он коммутирует со всеми
двенадцатью генераторами концевой калибровочной алгебры:
\begin{equation}
 [X_A,Y]=0,
 \qquad A=1,\ldots,12.
 \label{eq:v10-hypercharge-invariant-line}
\end{equation}
Кроме того, его представленная матрица на $21$-мерном концевом пространстве
имеет полный ранг $21$.

Поэтому вакуумная линия и шумовая линия гиперзаряда образуют
калибровочно-инвариантный двумерный подноситель
\begin{equation}
 \mathcal W_Y=\operatorname{span}\{|Y\rangle,|0\rangle\}
 \subset\mathcal K_{43}.
 \label{eq:v10-hypercharge-vacuum-subcarrier}
\end{equation}
Пусть $E_Y:\mathbb C^2\to\mathcal K_{43}$ задаётся столбцами
$E_Y=(|Y\rangle,|0\rangle)$. Тогда
\begin{equation}
 E_Y^*E_Y=I_2,
 \qquad P_Y=E_YE_Y^*,
 \qquad \rank P_Y=2.
 \label{eq:v10-k43-typed-isometry}
\end{equation}
Это не произвольный выбор двух координат: обе линии различены уже
существующими типами вакуума и центрального калибровочного направления.

\section{Совместимость со звёздным взаимодействием}

Старый гамильтониан столкновения содержит слагаемое
\begin{equation}
 H_Y=Y\otimes
 \bigl(|Y\rangle\langle0|+|0\rangle\langle Y|\bigr).
 \label{eq:v10-hypercharge-star-restriction}
\end{equation}
После ограничения на $\mathbb C^{21}\otimes\mathcal W_Y$ оно имеет ранг
\begin{equation}
 \rank H_Y=42.
 \label{eq:v10-hypercharge-star-rank}
\end{equation}
Следовательно, найденная подячейка не является молчащим наблюдателем: она
уже связана с каждой концевой компонентой через существующее физическое
направление гиперзаряда.

Тензорное поднятие
\begin{equation}
 E_Y^{\rm KMS}=E_Y\otimes I_6:
 \mathbb C^{12}\longrightarrow\mathcal K_{43}\otimes\mathbb C^6
 \label{eq:v10-k43-kms-lift}
\end{equation}
остаётся изометрией и имеет ранг $12$. Это доказывает совместимость типов с
KMS-множителем, но само по себе ещё не выводит условие детального равновесия
для нового спектрального оператора.

\section{Вложение взаимно обратного спектра}

На всей $43$-мерной ячейке определим оператор, тождественный на ортогональном
дополнении $\mathcal W_Y$:
\begin{equation}
 K_{43}(\zeta)
 =I_{43}+(e^{-\zeta}-1)|Y\rangle\langle Y|
 +(e^{\zeta}-1)|0\rangle\langle0|.
 \label{eq:v10-embedded-reciprocal-operator}
\end{equation}
Его сжатие точно воспроизводит оператор прошлого гейта,
\begin{equation}
 E_Y^*K_{43}(\zeta)E_Y
 =\operatorname{diag}(e^{-\zeta},e^{\zeta}),
 \qquad \det K_{43}(\zeta)=1.
 \label{eq:v10-embedded-reciprocal-compression}
\end{equation}
Обратный оператор получается заменой $\zeta\mapsto-\zeta$. Поэтому
самоэнергия входящей гиперзарядовой линии сохраняет точный ход
\begin{equation}
 \Sigma_Y(\zeta)
 =\langle Y|K_{43}^{-1}(\zeta)|Y\rangle=e^{\zeta},
 \qquad
 \frac{d\Sigma_Y}{d\zeta}=\Sigma_Y.
 \label{eq:v10-typed-hypercharge-self-energy}
\end{equation}

\section{Граница родительского происхождения}

Унаследованный локальный родитель вакуумной цепи равен
\begin{equation}
 h_{\rm cell}=I_{43}-|0\rangle\langle0|.
 \label{eq:v10-inherited-cell-parent}
\end{equation}
На найденной подячейке он даёт лишь
\begin{equation}
 E_Y^*h_{\rm cell}E_Y=\operatorname{diag}(1,0),
 \qquad
 \frac{d}{d\zeta}(E_Y^*h_{\rm cell}E_Y)=0.
 \label{eq:v10-inherited-cell-parent-zero-running}
\end{equation}
Таким образом, носитель, калибровочная линия, ненулевое взаимодействие и
KMS-типизация действительно унаследованы. Но коэффициенты $e^{\pm\zeta}$
пока внесены в новый оператор определением: старый родитель ячейки их не
порождает.

\begin{theorem}[Типизированный резервуар уже содержится в $\mathcal K_{43}$]
Пара из вакуума и гиперзарядового шумового направления задаёт каноническую
калибровочно-инвариантную изометрию $\mathbb C^2\hookrightarrow\mathcal
K_{43}$. Ограниченное звёздное взаимодействие имеет ранг $42$, KMS-поднятие
имеет ранг $12$, а вложенный взаимно обратный оператор сохраняет
$\Sigma_Y=e^{\zeta}$ и её ненулевой интенсивный ход.
\end{theorem}

\begin{nogocriterion}[Вложение носителя не выводит спектральный закон]
Из существования подпространства $\mathcal W_Y$ не следует зависимость
$K_{43}(\zeta)$. Унаследованный родитель $I-|0\rangle\langle0|$ не зависит
от $\zeta$ и имеет нулевую производную. Поэтому типизированное вложение
закрывает вопрос места действия, но не физическое происхождение взаимно
обратного спектра и абсолютного масштаба.
\end{nogocriterion}

\begin{protocol}[Статус типизированного вложения]\begin{enumerate}
\item изометрия $\mathbb C^2\hookrightarrow\mathcal K_{43}$:
\textbf{получена};
\item калибровочная инвариантность гиперзарядовой линии: \textbf{$12/12$};
\item ранг ограниченного звёздного взаимодействия: \textbf{$42$};
\item ранг KMS-поднятия: \textbf{$12$};
\item архитектура типизированного вложения: \textbf{$8/8$};
\item общий реестр происхождения: \textbf{$4/6$};
\item происхождение $e^{\pm\zeta}$ из старого родителя: \textbf{$0/1$};
\item общий родитель спектрального хода: \textbf{$0/1$};
\item абсолютный масштаб: \textbf{не выведен};
\item следующий гейт: \textbf{происхождение взаимно обратного
спектрального оператора из геометрического родителя роста}.
\end{enumerate}\end{protocol}

Машинный сертификат сохранён в
\path{s2t_v10_inflow_spectral_self_energy_k43_typed_embedding_gate_results.json}.